%----------------------------------------------------------------------------------------
%    PACKAGES AND THEMES
%----------------------------------------------------------------------------------------

\documentclass[aspectratio=169,xcolor=dvipsnames]{beamer}

\input{slides/utils/preamble}
\input{slides/utils/custom-commands}

\usetikzlibrary{decorations.pathreplacing}

%----------------------------------------------------------------------------------------
%    CONTAGEM ANIMADA DE PARÂMETROS DA VGG16
%    \vstep indica o passo atual: cada elemento do diagrama e cada bloco da pilha
%    é "futuro" (apagado/oculto), "atual" (destacado) ou "contado" (normal).
%----------------------------------------------------------------------------------------
\newcommand{\vstep}{14}
\newcommand{\vgpitch}{0.27}
\tikzset{
    % elemento do diagrama: {passo}{cor}
    vgsty/.code 2 args={%
        \ifnum#1>\vstep\space \pgfkeysalso{fill=black!6, draw=black!25, thin}%
        \else\ifnum#1=\vstep\space \pgfkeysalso{fill=#2, draw=black, line width=1.3pt}%
        \else \pgfkeysalso{fill=#2, draw=black!50, thin}\fi\fi},
    % rótulo do diagrama: passo
    vglab/.code={%
        \ifnum#1>\vstep\space \pgfkeysalso{text=black!35}%
        \else\ifnum#1=\vstep\space \pgfkeysalso{text=black, font=\tiny\bfseries}%
        \else \pgfkeysalso{text=black!75}\fi\fi},
    % bloco da pilha: {passo}{cor de fundo}{cor do texto}
    vgchip/.code n args={3}{%
        \ifnum#1>\vstep\space \pgfkeysalso{opacity=0}%
        \else\ifnum#1=\vstep\space \pgfkeysalso{fill=#2, text=#3, draw=black, line width=1.2pt}%
        \else \pgfkeysalso{fill=#2, text=#3, draw=black!50, thin}\fi\fi},
    vgchipbase/.style={minimum width=1.9cm, minimum height=0.22cm, inner sep=1pt, rounded corners=1pt, font=\tiny},
}
\newcommand{\chipconv}[1]{conv $3{\times}3$, #1}
\newcommand{\chipfc}[1]{fc, #1}

% Pilha de camadas (de baixo para cima) com centro do primeiro bloco em (#1,#2)
\newcommand{\vggstack}[2]{%
    \foreach \i/\s/\txt/\col/\tcol in {
        0/0/entrada/black!55/white,
        1/1/\chipconv{64}/BrickRed!75/white,
        2/1/\chipconv{64}/BrickRed!75/white,
        3/2/pooling/PastelGreen/black,
        4/3/\chipconv{128}/BrickRed!75/white,
        5/3/\chipconv{128}/BrickRed!75/white,
        6/4/pooling/PastelGreen/black,
        7/5/\chipconv{256}/BrickRed!75/white,
        8/5/\chipconv{256}/BrickRed!75/white,
        9/5/\chipconv{256}/BrickRed!75/white,
        10/6/pooling/PastelGreen/black,
        11/7/\chipconv{512}/BrickRed!75/white,
        12/7/\chipconv{512}/BrickRed!75/white,
        13/7/\chipconv{512}/BrickRed!75/white,
        14/8/pooling/PastelGreen/black,
        15/9/\chipconv{512}/BrickRed!75/white,
        16/9/\chipconv{512}/BrickRed!75/white,
        17/9/\chipconv{512}/BrickRed!75/white,
        18/10/pooling/PastelGreen/black,
        19/11/\chipfc{4096}/PastelBlue2!80/white,
        20/11/\chipfc{4096}/PastelBlue2!80/white,
        21/12/\chipfc{1000}/PastelBlue2!80/white,
        22/13/softmax/Goldenrod!70/black}{
        \node[vgchipbase, vgchip={\s}{\col}{\tcol}] at (#1, #2 + \i*\vgpitch) {\txt};
    }%
}

% Texto da caixa de total acumulado
\newcommand{\vggtotaltext}{%
    \ifcase\vstep
        \textbf{Total de parâmetros treináveis: 0}\\ a entrada não possui parâmetros\or
        \textbf{Total de parâmetros treináveis: 38.592}\\ \mbox{$3{\cdot}3{\cdot}3{\cdot}64 + 3{\cdot}3{\cdot}64{\cdot}64$}\\ \mbox{$= 1.728 + 36.864 = 38.592$}\or
        \textbf{Total de parâmetros treináveis: 38.592}\\ \textit{max pooling} não possui parâmetros (\mbox{$+0$})\or
        \textbf{Total de parâmetros treináveis: 259.776}\\ \mbox{$3{\cdot}3{\cdot}64{\cdot}128 + 3{\cdot}3{\cdot}128{\cdot}128$}\\ \mbox{$= 73.728 + 147.456 = 221.184$}\or
        \textbf{Total de parâmetros treináveis: 259.776}\\ \textit{max pooling} não possui parâmetros (\mbox{$+0$})\or
        \textbf{Total de parâmetros treináveis: 1.734.336}\\ \mbox{$3{\cdot}3{\cdot}128{\cdot}256 + 3{\cdot}3{\cdot}256{\cdot}256 + 3{\cdot}3{\cdot}256{\cdot}256$}\\ \mbox{$= 294.912 + 589.824 + 589.824 = 1.474.560$}\or
        \textbf{Total de parâmetros treináveis: 1.734.336}\\ \textit{max pooling} não possui parâmetros (\mbox{$+0$})\or
        \textbf{Total de parâmetros treináveis: 7.632.576}\\ \mbox{$3{\cdot}3{\cdot}256{\cdot}512 + 3{\cdot}3{\cdot}512{\cdot}512 + 3{\cdot}3{\cdot}512{\cdot}512$}\\ \mbox{$= 1.179.648 + 2.359.296 + 2.359.296 = 5.898.240$}\or
        \textbf{Total de parâmetros treináveis: 7.632.576}\\ \textit{max pooling} não possui parâmetros (\mbox{$+0$})\or
        \textbf{Total de parâmetros treináveis: 14.710.464}\\ \mbox{$3 \times (3{\cdot}3{\cdot}512{\cdot}512)$}\\ \mbox{$= 3 \times 2.359.296 = 7.077.888$}\or
        \textbf{Total de parâmetros treináveis: 14.710.464}\\ \textit{max pooling} não possui parâmetros (\mbox{$+0$})\or
        \textbf{Total de parâmetros treináveis: 134.248.128}\\ \mbox{$7{\cdot}7{\cdot}512{\cdot}4096 + 4096{\cdot}4096$}\\ \mbox{$= 102.760.448 + 16.777.216 = 119.537.664$}\or
        \textbf{Total de parâmetros treináveis: 138.344.128}\\ \mbox{$4096{\cdot}1000 = 4.096.000$}\or
        \textbf{Total de parâmetros treináveis: 138.344.128}\\ softmax não possui parâmetros (\mbox{$+0$})\or
        \textbf{Total de parâmetros treináveis: 138.344.128}\\ Observação: não contamos os \textit{bias}\\ (mais \mbox{$13.416$}, totalizando \mbox{$138.357.544$})\fi
}

% Figura completa: caixa de total, diagrama com destaque e pilha à direita
\newcommand{\vggcountpic}{%
\begin{tikzpicture}[every node/.style={font=\tiny}]
    % caixa de total acumulado
    \node[rounded corners=4pt, fill=black!70, text=white, align=center, text width=10.3cm,
          minimum height=1.6cm, inner sep=4pt, font=\small] at (5.15,5.45) {\vggtotaltext};

    % diagrama da arquitetura
    \begin{scope}[shift={(0,1.8)}, scale=0.75]
        \def\bw{0.55}
        % entrada
        \draw[vgsty={0}{black!20}] (0,0) rectangle ++(\bw,2.6);
        \node[vglab=0, align=center] at (\bw/2,-0.35) {$224{\times}224$\\$\times 3$};
        % bloco 1
        \draw[vgsty={1}{BrickRed!25}] (1.0,0) rectangle ++(\bw,2.6);
        \draw[vgsty={1}{BrickRed!25}] (1.6,0) rectangle ++(\bw,2.6);
        \node[vglab=1, align=center] at (1.55,-0.35) {$224{\times}224$\\$\times 64$};
        % pooling 1
        \draw[vgsty={2}{PastelGreen}] (2.3,0.3) rectangle ++(0.22,2.0);
        % bloco 2
        \draw[vgsty={3}{BrickRed!40}] (2.75,0.3) rectangle ++(\bw,2.0);
        \draw[vgsty={3}{BrickRed!40}] (3.35,0.3) rectangle ++(\bw,2.0);
        \node[vglab=3, align=center] at (3.3,-0.35) {$112{\times}112$\\$\times 128$};
        % pooling 2
        \draw[vgsty={4}{PastelGreen}] (4.05,0.55) rectangle ++(0.22,1.5);
        % bloco 3
        \draw[vgsty={5}{BrickRed!55}] (4.5,0.55) rectangle ++(\bw,1.5);
        \draw[vgsty={5}{BrickRed!55}] (5.1,0.55) rectangle ++(\bw,1.5);
        \draw[vgsty={5}{BrickRed!55}] (5.7,0.55) rectangle ++(\bw,1.5);
        \node[vglab=5, align=center] at (5.35,-0.35) {$56{\times}56$\\$\times 256$};
        % pooling 3
        \draw[vgsty={6}{PastelGreen}] (6.4,0.75) rectangle ++(0.22,1.1);
        % bloco 4
        \draw[vgsty={7}{BrickRed!70}] (6.85,0.75) rectangle ++(\bw,1.1);
        \draw[vgsty={7}{BrickRed!70}] (7.45,0.75) rectangle ++(\bw,1.1);
        \draw[vgsty={7}{BrickRed!70}] (8.05,0.75) rectangle ++(\bw,1.1);
        \node[vglab=7, align=center] at (7.7,-0.35) {$28{\times}28$\\$\times 512$};
        % pooling 4
        \draw[vgsty={8}{PastelGreen}] (8.75,0.95) rectangle ++(0.22,0.7);
        % bloco 5
        \draw[vgsty={9}{BrickRed!85}] (9.2,0.95) rectangle ++(\bw,0.7);
        \draw[vgsty={9}{BrickRed!85}] (9.8,0.95) rectangle ++(\bw,0.7);
        \draw[vgsty={9}{BrickRed!85}] (10.4,0.95) rectangle ++(\bw,0.7);
        \node[vglab=9, align=center] at (10.05,-0.35) {$14{\times}14$\\$\times 512$};
        % pooling 5 (saída achatada)
        \draw[vgsty={10}{PastelGreen}] (11.15,1.15) rectangle ++(0.4,0.3);
        \node[vglab=10, align=center] at (11.35,-0.35) {$7{\times}7$\\$\times 512$};
        % densas
        \draw[vgsty={11}{PastelBlue2!80}] (11.95,1.2) rectangle ++(0.25,0.2);
        \draw[vgsty={11}{PastelBlue2!80}] (12.35,1.2) rectangle ++(0.25,0.2);
        \node[vglab=11, align=center] at (12.28,-0.35) {fc\\4096};
        \draw[vgsty={12}{PastelBlue2!80}] (12.95,1.2) rectangle ++(0.25,0.2);
        \node[vglab=12, align=center] at (13.07,-0.35) {fc\\1000};
        % softmax
        \draw[vgsty={13}{Goldenrod!70}] (13.6,1.2) rectangle ++(0.25,0.2);
        \node[vglab=13, align=center] at (13.72,-0.35) {soft\\max};
        % legenda
        \node[anchor=west, text=black!75] at (0,-1.2) {%
            \tikz \fill[BrickRed!50] (0,0) rectangle (0.25,0.15); ~conv $3\times 3$ + ReLU \quad
            \tikz \fill[PastelGreen] (0,0) rectangle (0.25,0.15); ~\textit{max pooling} \quad
            \tikz \fill[PastelBlue2!80] (0,0) rectangle (0.25,0.15); ~totalmente conectada \quad
            \tikz \fill[Goldenrod!70] (0,0) rectangle (0.25,0.15); ~softmax};
    \end{scope}

    % pilha de camadas à direita
    \vggstack{11.6}{0}
\end{tikzpicture}}

%----------------------------------------------------------------------------------------
%    PILHA PARA TRANSFERÊNCIA DE APRENDIZADO
%    \vgtransferstack{x}{y}{modo}: modo 1 = convoluções congeladas + cabeça nova,
%    modo 2 = fine-tuning da rede inteira (cabeça nova continua destacada).
%    Cada bloco carrega uma parte: 0 = entrada, 1 = extrator convolucional, 2 = cabeça.
%----------------------------------------------------------------------------------------
\tikzset{
    % bloco da pilha de transferência: {parte}{modo}
    vgtchip/.code 2 args={%
        \ifnum#2=1
            \ifnum#1=2 \pgfkeysalso{draw=black, line width=1.2pt}%
            \else\ifnum#1=1 \pgfkeysalso{fill opacity=0.3, text opacity=1, text=black!60, draw=black!45, thin, dashed}%
            \else \pgfkeysalso{draw=black!50, thin}\fi\fi
        \else
            \ifnum#1=2 \pgfkeysalso{draw=black, line width=1.2pt}%
            \else \pgfkeysalso{draw=black!50, thin}\fi
        \fi},
}
\newcommand{\vgtransferstack}[3]{%
    \foreach \i/\p/\txt/\col/\tcol in {
        0/0/entrada/black!55/white,
        1/1/\chipconv{64}/BrickRed!75/white,
        2/1/\chipconv{64}/BrickRed!75/white,
        3/1/pooling/PastelGreen/black,
        4/1/\chipconv{128}/BrickRed!75/white,
        5/1/\chipconv{128}/BrickRed!75/white,
        6/1/pooling/PastelGreen/black,
        7/1/\chipconv{256}/BrickRed!75/white,
        8/1/\chipconv{256}/BrickRed!75/white,
        9/1/\chipconv{256}/BrickRed!75/white,
        10/1/pooling/PastelGreen/black,
        11/1/\chipconv{512}/BrickRed!75/white,
        12/1/\chipconv{512}/BrickRed!75/white,
        13/1/\chipconv{512}/BrickRed!75/white,
        14/1/pooling/PastelGreen/black,
        15/1/\chipconv{512}/BrickRed!75/white,
        16/1/\chipconv{512}/BrickRed!75/white,
        17/1/\chipconv{512}/BrickRed!75/white,
        18/1/pooling/PastelGreen/black,
        19/2/\chipfc{4096}/PastelBlue2!80/white,
        20/2/\chipfc{4096}/PastelBlue2!80/white,
        21/2/\chipfc{4}/PastelBlue2!80/white,
        22/2/sigmoide/Goldenrod!70/black}{
        \node[vgchipbase, fill=\col, text=\tcol, vgtchip={\p}{#3}] at (#1, #2 + \i*\vgpitch) {\txt};
    }%
}

%----------------------------------------------------------------------------------------
%    TITLE PAGE
%----------------------------------------------------------------------------------------

\title{Deep Learning}
\subtitle{Aplicações de CNN em Visão Computacional}

\author{Lucas Silveira Kupssinskü}

\institute
{
    Machine Learning Theory and Applications Laboratory \\
    PUCRS
}
\date{\mespor de \the\year}

%----------------------------------------------------------------------------------------
%    PRESENTATION SLIDES
%----------------------------------------------------------------------------------------

\begin{document}

\begin{frame}
    \titlepage
\end{frame}

\begin{frame}{Visão Geral}
    \tableofcontents
\end{frame}

%========================================================================
\section{VGG}
%========================================================================

\begin{frame}{VGG}
\begin{itemize}
    \item Arquitetura de CNN proposta em 2014, publicada no ICLR 2015\footfullcite{simonyan2014vgg}.
    \item Continua sendo usada em muitos cenários, especialmente como \textit{backbone} pré-treinado.
    \item No ILSVRC-2014 (\textit{ImageNet Large Scale Visual Recognition Challenge}) ficou em \textbf{1º na localização} e \textbf{2º na classificação} (vencida pelo GoogLeNet).
    \item Taxa de erro no ImageNet (\textit{ensemble multi-crop} reportado no artigo):
    \begin{itemize}
        \item \textit{Top-5}: $6.8\%$.
        \item \textit{Top-1}: $23.7\%$.
        \item VGG16 \textit{single-model} fica em torno de $7.3\%$ \textit{top-5} e $24.7\%$ \textit{top-1}.
    \end{itemize}
    \item $\sim 140$M parâmetros treináveis.
\end{itemize}
\end{frame}

\begin{frame}{Definição do Problema}
\begin{itemize}
    \item ImageNet contém imagens RGB na resolução $224\times 224 \times 3$.
    \begin{itemize}
        \item As entradas são tensores $(224, 224, 3)$.
    \end{itemize}
    \item O objetivo é classificar cada imagem em uma das $1000$ classes pré-definidas.
    \begin{itemize}
        \item Função de custo: entropia cruzada com os \textit{logits}.
    \end{itemize}
    \item Precisamos transformar um tensor $(224, 224, 3)$ em um vetor $(1, 1, 1000)$.
\end{itemize}
\end{frame}

\begin{frame}{Blocos de Construção}
\begin{itemize}
    \item A VGG combina três tipos de operações:
    \begin{itemize}
        \item \textbf{Convolução + ReLU}: aumenta o número de canais e extrai características espaciais.
        \item \textbf{\textit{Pooling}}: subamostra a dimensão espacial (largura e altura).
        \item \textbf{Camadas densas}: realizam a classificação final, partindo do \textit{feature map} achatado.
    \end{itemize}
    \item Toda a parte convolucional pode ser vista como um \textbf{extrator de \textit{features}}.
    \item As camadas densas no fim convertem essas \textit{features} em probabilidades de classe.
\end{itemize}
\end{frame}

\begin{frame}{VGG16}
\begin{itemize}
    \item VGG16 é a implementação mais comum, equivalente à variante ``D'' descrita no artigo.
    \item Possui 16 camadas com parâmetros treináveis: \textcolor{BrickRed}{\textbf{13 convolucionais}} e \textcolor{PastelBlue2}{\textbf{3 densas}}.
    \item Todas as convoluções são $3\times 3$ com \textit{stride} 1 e \textit{padding} 1 (mantêm a dimensão espacial).
    \item Cada \textit{max pooling} usa janela $2\times 2$ com \textit{stride} 2 (reduz pela metade a largura e altura).
\end{itemize}
\end{frame}

\begin{frame}{Arquitetura da VGG16}
\centering
\begin{tikzpicture}[every node/.style={font=\scriptsize}, scale=0.95, xscale=1.04]
    % Simple cuboid-like rectangles suggest tensor volumes (same layout as the parameter-counting frames).
    \def\bw{0.55} % box width

    % Input
    \draw[fill=black!20, draw=black!50] (0,0) rectangle ++(\bw,2.6);
    \node[align=center] at (\bw/2,-0.35) {$224{\times}224$\\$\times 3$};

    % Block 1: conv 64 x2
    \draw[fill=BrickRed!25] (1.0,0) rectangle ++(\bw,2.6);
    \draw[fill=BrickRed!25] (1.6,0) rectangle ++(\bw,2.6);
    \node[align=center] at (1.55,-0.35) {$224{\times}224$\\$\times 64$};
    % Pool 1
    \draw[fill=PastelGreen] (2.3,0.3) rectangle ++(0.22,2.0);

    % Block 2: conv 128 x2
    \draw[fill=BrickRed!40] (2.75,0.3) rectangle ++(\bw,2.0);
    \draw[fill=BrickRed!40] (3.35,0.3) rectangle ++(\bw,2.0);
    \node[align=center] at (3.3,-0.35) {$112{\times}112$\\$\times 128$};
    % Pool 2
    \draw[fill=PastelGreen] (4.05,0.55) rectangle ++(0.22,1.5);

    % Block 3: conv 256 x3
    \draw[fill=BrickRed!55] (4.5,0.55) rectangle ++(\bw,1.5);
    \draw[fill=BrickRed!55] (5.1,0.55) rectangle ++(\bw,1.5);
    \draw[fill=BrickRed!55] (5.7,0.55) rectangle ++(\bw,1.5);
    \node[align=center] at (5.35,-0.35) {$56{\times}56$\\$\times 256$};
    % Pool 3
    \draw[fill=PastelGreen] (6.4,0.75) rectangle ++(0.22,1.1);

    % Block 4: conv 512 x3
    \draw[fill=BrickRed!70] (6.85,0.75) rectangle ++(\bw,1.1);
    \draw[fill=BrickRed!70] (7.45,0.75) rectangle ++(\bw,1.1);
    \draw[fill=BrickRed!70] (8.05,0.75) rectangle ++(\bw,1.1);
    \node[align=center] at (7.7,-0.35) {$28{\times}28$\\$\times 512$};
    % Pool 4
    \draw[fill=PastelGreen] (8.75,0.95) rectangle ++(0.22,0.7);

    % Block 5: conv 512 x3
    \draw[fill=BrickRed!85] (9.2,0.95) rectangle ++(\bw,0.7);
    \draw[fill=BrickRed!85] (9.8,0.95) rectangle ++(\bw,0.7);
    \draw[fill=BrickRed!85] (10.4,0.95) rectangle ++(\bw,0.7);
    \node[align=center] at (10.05,-0.35) {$14{\times}14$\\$\times 512$};
    % Pool 5 (flattened output)
    \draw[fill=PastelGreen] (11.15,1.15) rectangle ++(0.4,0.3);
    \node[align=center] at (11.35,-0.35) {$7{\times}7$\\$\times 512$};

    % Dense layers
    \draw[fill=PastelBlue2!80] (11.95,1.2) rectangle ++(0.25,0.2);
    \draw[fill=PastelBlue2!80] (12.35,1.2) rectangle ++(0.25,0.2);
    \node[align=center] at (12.28,-0.35) {fc\\4096};
    \draw[fill=PastelBlue2!80] (12.95,1.2) rectangle ++(0.25,0.2);
    \node[align=center] at (13.07,-0.35) {fc\\1000};

    % Output (softmax)
    \draw[fill=Goldenrod!70] (13.6,1.2) rectangle ++(0.25,0.2);
    \node[align=center] at (13.72,-0.35) {soft\\max};

    % Legend
    \node[anchor=west, font=\tiny] at (0,3.0) {\tikz \fill[BrickRed!50] (0,0) rectangle (0.25,0.18); ~conv $3\times 3$ + ReLU \quad
        \tikz \fill[PastelGreen] (0,0) rectangle (0.25,0.18); ~\textit{max pooling} \quad
        \tikz \fill[PastelBlue2!80] (0,0) rectangle (0.25,0.18); ~totalmente conectada \quad
        \tikz \fill[Goldenrod!70] (0,0) rectangle (0.25,0.18); ~softmax};
\end{tikzpicture}

\vspace{0.3em}
\footnotesize Cada \textit{max pooling} reduz a dimensão espacial pela metade e a convolução seguinte dobra (ou mantém) o número de canais.
\end{frame}

\begin{frame}{Contando Parâmetros: Regras}
\begin{itemize}
    \item Convolução com \textit{kernel} $K_h \times K_w$, $C_{in}$ canais de entrada e $C_{out}$ filtros: $K_h \cdot K_w \cdot C_{in} \cdot C_{out}$ pesos.
    \item Camada densa com $n_{in}$ entradas e $n_{out}$ saídas: $n_{in} \cdot n_{out}$ pesos.
    \item \textit{Max pooling}, ReLU e softmax não possuem parâmetros treináveis.
    \item Na VGG16 toda convolução é $3 \times 3$, então cada uma custa $9 \cdot C_{in} \cdot C_{out}$.
\end{itemize}
\begin{block}{Observação}
Ignoramos o \textit{bias} para simplificar a contagem. Na prática, cada filtro e cada neurônio denso têm um \textit{bias} associado.
\end{block}
\end{frame}

\begin{frame}{Contando Parâmetros da VGG16}
\centering
\only<1>{\renewcommand{\vstep}{0}\vggcountpic}%
\only<2>{\renewcommand{\vstep}{1}\vggcountpic}%
\only<3>{\renewcommand{\vstep}{2}\vggcountpic}%
\only<4>{\renewcommand{\vstep}{3}\vggcountpic}%
\only<5>{\renewcommand{\vstep}{4}\vggcountpic}%
\only<6>{\renewcommand{\vstep}{5}\vggcountpic}%
\only<7>{\renewcommand{\vstep}{6}\vggcountpic}%
\only<8>{\renewcommand{\vstep}{7}\vggcountpic}%
\only<9>{\renewcommand{\vstep}{8}\vggcountpic}%
\only<10>{\renewcommand{\vstep}{9}\vggcountpic}%
\only<11>{\renewcommand{\vstep}{10}\vggcountpic}%
\only<12>{\renewcommand{\vstep}{11}\vggcountpic}%
\only<13>{\renewcommand{\vstep}{12}\vggcountpic}%
\only<14>{\renewcommand{\vstep}{13}\vggcountpic}%
\only<15>{\renewcommand{\vstep}{14}\vggcountpic}%
\end{frame}

\begin{frame}{Onde Estão os Parâmetros?}
\begin{columns}
    \begin{column}{.58\linewidth}
        \begin{itemize}
            \item A maioria dos parâmetros treináveis está nas \textbf{camadas densas}.
            \begin{itemize}
                \item Convoluções: $\sim 14$M.
                \item Densas: $\sim 124$M.
            \end{itemize}
            \item Convoluções compartilham pesos, então usam ordens de magnitude menos parâmetros que camadas totalmente conectadas equivalentes.
            \item Por isso é comum trocar as densas por uma cabeça menor quando se reaproveita a VGG em outras tarefas.
        \end{itemize}
    \end{column}
    \begin{column}{.42\linewidth}
        \centering
        \renewcommand{\vstep}{14}%
        \begin{tikzpicture}[every node/.style={font=\tiny}]
            \vggstack{0}{0}
            \draw[decorate, decoration={brace, amplitude=6pt}, BrickRed, thick]
                (-1.05, 1*\vgpitch - 0.12) -- (-1.05, 18*\vgpitch + 0.12)
                node[midway, left=8pt, BrickRed, font=\Large\bfseries] {$\sim 14$M};
            \draw[decorate, decoration={brace, amplitude=6pt}, PastelBlue2, thick]
                (-1.05, 19*\vgpitch - 0.12) -- (-1.05, 21*\vgpitch + 0.12)
                node[midway, left=8pt, PastelBlue2, font=\Large\bfseries] {$\sim 124$M};
        \end{tikzpicture}
    \end{column}
\end{columns}
\end{frame}

%========================================================================
\section{Localização}
%========================================================================

\begin{frame}{Classificação vs Classificação + Localização}
\begin{itemize}
    \item Além de dizer ``há um gato na imagem'', queremos saber \textbf{onde} esse gato está.
    \item A saída adicional é uma \textit{bounding box} delimitando o objeto.
\end{itemize}
\vspace{0.4em}
\centering
\begin{tikzpicture}[every node/.style={font=\scriptsize}]
    % Classification (image 4.2cm x 2.8cm, 3:2)
    \node[inner sep=0, anchor=north west] at (0,0)
      {\includegraphics[width=4.2cm]{slides/img/02-cnn/loc_cat_snow.jpg}};
    \node[anchor=south] at (2.1,0.05) {\small Classificação};
    \node[anchor=north] at (2.1,-2.85) {saída: classe \textbf{gato}};

    % Classification + Localization (bbox in normalized coords)
    \node[inner sep=0, anchor=north west] at (5.4,0)
      {\includegraphics[width=4.2cm]{slides/img/02-cnn/loc_cat_snow.jpg}};
    \node[anchor=south] at (7.5,0.05) {\small Classificação + Localização};
    \begin{scope}[shift={(5.4,0)}, x=4.2cm, y=-2.8cm]
        \draw[BrickRed, very thick] (0.21,0.10) rectangle (0.97,0.93);
        \node[BrickRed, anchor=north east, fill=white, inner sep=1pt, font=\scriptsize\bfseries] at (0.97,0.10) {gato};
    \end{scope}
    \node[anchor=north] at (7.5,-2.85) {saída: classe \textbf{gato} $+$ \textit{bounding box}};
\end{tikzpicture}\\[2pt]
{\tiny Foto: Von.grzanka, \textit{Wikimedia Commons} (CC BY-SA 3.0).}
\end{frame}

\begin{frame}{Bounding Box}
\begin{itemize}
    \item Para descrever uma \textit{bounding box} precisamos de 4 valores:
    \begin{itemize}
        \item $(x, y)$: coordenadas de um ponto de referência (canto superior esquerdo ou centro).
        \item $h$: altura da caixa.
        \item $w$: largura da caixa.
    \end{itemize}
    \item Podemos modelar esses 4 valores como um problema de \textbf{regressão multivariada}.
    \item Convenção comum: o canto superior esquerdo da imagem é $(0,0)$ e o canto inferior direito é $(1,1)$.
\end{itemize}
\end{frame}

\begin{frame}{Exemplo de Bounding Box}
\begin{columns}
    \begin{column}{.5\linewidth}
        \begin{itemize}
            \item Alvo $\mathbf{y} = (x, y, h, w)$.
            \item Para o exemplo ao lado:
            \begin{itemize}
                \item $(x, y) = (0.21,\, 0.10)$.
                \item $h = 0.83$.
                \item $w = 0.76$.
            \end{itemize}
            \item Logo, $\mathbf{y} = (0.21,\, 0.10,\, 0.83,\, 0.76)$.
        \end{itemize}
    \end{column}
    \begin{column}{.5\linewidth}
        \centering
        \begin{tikzpicture}[every node/.style={font=\scriptsize}]
            % Image 6cm x 4cm (3:2); scope maps normalized coords, origin top-left
            \node[inner sep=0, anchor=north west] at (0,0)
              {\includegraphics[width=6cm]{slides/img/02-cnn/loc_cat_snow.jpg}};
            \draw[thin, black!60] (0,0) rectangle (6,-4);
            \begin{scope}[x=6cm, y=-4cm]
                \node[anchor=south west, font=\tiny, inner sep=1pt] at (0,0) {$(0, 0)$};
                \node[anchor=north east, font=\tiny, inner sep=1pt] at (1,1) {$(1, 1)$};
                % Bounding box: x=0.21, y=0.10, w=0.76, h=0.83
                \draw[BrickRed, very thick] (0.21,0.10) rectangle (0.97,0.93);
                \fill[BrickRed] (0.21,0.10) circle (1.2pt);
                \node[anchor=south east, BrickRed, font=\tiny, fill=white, inner sep=1pt] at (0.21,0.10) {$(0.21,\, 0.10)$};
                \draw[<->, BrickRed, thick] (0.15,0.10) -- (0.15,0.93) node[midway, left, font=\tiny, fill=white, inner sep=1pt] {$h=0.83$};
                \draw[<->, BrickRed, thick] (0.21,0.97) -- (0.97,0.97) node[midway, below, font=\tiny, fill=white, inner sep=1pt] {$w=0.76$};
            \end{scope}
        \end{tikzpicture}\\[4pt]
        {\tiny Foto: Von.grzanka, \textit{Wikimedia Commons} (CC BY-SA 3.0).}
    \end{column}
\end{columns}
\end{frame}

\begin{frame}{Arquitetura para Localização}
\begin{itemize}
    \item Para regressão dos 4 valores da \textit{bounding box}, a saída precisa de \textbf{4 unidades sigmoide} (em vez de \textit{softmax} para 1000 classes).
    \begin{itemize}
        \item A sigmoide garante saídas em $[0, 1]$, compatíveis com coordenadas normalizadas.
        \item Hipótese implícita: a caixa cabe inteiramente dentro da imagem, com $(x, y) \in [0,1]^2$ e $h, w \in [0,1]$.
        \item Caixas que estouram a borda exigem outra parametrização (por exemplo, \textit{logits} sem ativação ou regressão de \textit{offsets} em relação a uma \textit{anchor}).
    \end{itemize}
    \item A função de custo passa a ser o \textbf{erro quadrático médio} (MSE):
\end{itemize}
\begin{equation*}
\mathcal{L} = \frac{1}{2N}\sum_{i=1}^{N}\Bigl[\bigl(y_0^{(i)} - \hat{y}_0^{(i)}\bigr)^2 + \bigl(y_1^{(i)} - \hat{y}_1^{(i)}\bigr)^2 + \bigl(y_2^{(i)} - \hat{y}_2^{(i)}\bigr)^2 + \bigl(y_3^{(i)} - \hat{y}_3^{(i)}\bigr)^2\Bigr]
\end{equation*}
\begin{itemize}
    \item As convoluções continuam funcionando como extrator de \textit{features}.
\end{itemize}
\end{frame}

\begin{frame}{Reaproveitando uma Rede Pré-treinada}
\begin{columns}[T]
    \begin{column}{.55\linewidth}
        \begin{itemize}
            \item<1-> Imagine que já treinamos a rede para \textbf{classificar} as imagens.
            \begin{itemize}
                \item Os parâmetros das convoluções já estão ajustados para extrair \textit{features} úteis.
            \end{itemize}
            \item<2-> Estratégia comum: \textbf{congelar} os pesos da convolução e treinar apenas uma \textbf{cabeça densa nova} para a nova tarefa\footfullcite{sermanet2013overfeat}.
            \begin{itemize}
                \item Reduz drasticamente o tempo de treinamento.
                \item Funciona bem quando o conjunto de dados novo é pequeno.
            \end{itemize}
            \item<3-> Alternativa: \textit{fine-tuning} da rede inteira com \textit{learning rate} baixa.
        \end{itemize}
    \end{column}
    \begin{column}{.45\linewidth}
        \centering
        \renewcommand{\vgpitch}{0.245}%
        \begin{tikzpicture}[every node/.style={font=\tiny}]
            \only<1>{%
                \renewcommand{\vstep}{14}%
                \vggstack{0}{0}
                \draw[decorate, decoration={brace, amplitude=6pt}, black!70, thick]
                    (-1.05, 1*\vgpitch - 0.1) -- (-1.05, 22*\vgpitch + 0.1)
                    node[midway, left=8pt, black!70, font=\small\bfseries, align=center] {rede treinada\\para classificar\\1000 classes};
            }%
            \only<2>{%
                \vgtransferstack{0}{0}{1}
                \draw[decorate, decoration={brace, amplitude=6pt}, black!60, thick]
                    (-1.05, 1*\vgpitch - 0.1) -- (-1.05, 18*\vgpitch + 0.1)
                    node[midway, left=8pt, black!60, font=\small\bfseries, align=center] {congelado\\(reaproveitado)};
                \draw[decorate, decoration={brace, amplitude=6pt}, PastelBlue2, thick]
                    (-1.05, 19*\vgpitch - 0.1) -- (-1.05, 22*\vgpitch + 0.1)
                    node[midway, left=8pt, PastelBlue2, font=\small\bfseries, align=center] {cabeça nova\\(treinada)};
            }%
            \only<3>{%
                \vgtransferstack{0}{0}{2}
                \draw[decorate, decoration={brace, amplitude=6pt}, BrickRed, thick]
                    (-1.05, 1*\vgpitch - 0.1) -- (-1.05, 22*\vgpitch + 0.1)
                    node[midway, left=8pt, BrickRed, font=\small\bfseries, align=center] {rede inteira\\treinada\\(\textit{learning rate}\\baixa)};
                \draw[decorate, decoration={brace, amplitude=6pt, mirror}, PastelBlue2, thick]
                    (1.05, 19*\vgpitch - 0.1) -- (1.05, 22*\vgpitch + 0.1)
                    node[midway, right=6pt, PastelBlue2, font=\small\bfseries, align=center] {cabeça\\nova};
            }%
        \end{tikzpicture}
    \end{column}
\end{columns}
\end{frame}

\begin{frame}{Métricas de Localização: \textit{Intersection over Union}}
\begin{columns}
    \begin{column}{.55\linewidth}
        \begin{itemize}
            \item Computa a interseção entre a \textit{bounding box} prevista e o \textit{ground truth}, sobre a união das duas áreas.
            \begin{equation*}
                \mathrm{IoU} = \frac{A \cap B}{A \cup B}
            \end{equation*}
            \item Valor ideal: $\mathrm{IoU} = 1$ (sobreposição perfeita).
            \item Ignora o tamanho absoluto das caixas: só importa a sobreposição relativa.
        \end{itemize}
    \end{column}
    \begin{column}{.45\linewidth}
        \centering
        \begin{tikzpicture}
            % Ground truth
            \draw[thick, BrickRed, fill=BrickRed, fill opacity=0.25] (0,0.6) rectangle (2.6,2.6);
            \node[BrickRed, font=\scriptsize] at (0.5,2.85) {ground truth $A$};
            % Prediction
            \draw[thick, PastelBlue2, fill=PastelBlue2, fill opacity=0.25] (1.4,0) rectangle (3.6,2.0);
            \node[PastelBlue2, font=\scriptsize] at (3.0,-0.3) {predição $B$};
            % Intersection emphasis
            \draw[ultra thick, dashed, black] (1.4,0.6) rectangle (2.6,2.0);
            \node[font=\scriptsize] at (2.0,1.3) {$A \cap B$};
        \end{tikzpicture}
    \end{column}
\end{columns}
\end{frame}

\begin{frame}{Métricas de Localização: \textit{Dice Score}}
\begin{columns}
    \begin{column}{.55\linewidth}
        \begin{itemize}
            \item Computa o dobro da interseção sobre a soma das áreas.
            \begin{equation*}
                \mathrm{Dice} = \frac{2 \cdot (A \cap B)}{\mathrm{Area}(A) + \mathrm{Area}(B)}
            \end{equation*}
            \item Valor ideal: $\mathrm{Dice} = 1$.
            \item Penaliza diferenças de tamanho menos que o IoU.
        \end{itemize}
    \end{column}
    \begin{column}{.45\linewidth}
        \centering
        \begin{tikzpicture}
            \draw[thick, BrickRed, fill=BrickRed, fill opacity=0.25] (0,0.6) rectangle (2.6,2.6);
            \node[BrickRed, font=\scriptsize] at (0.5,2.85) {$A$};
            \draw[thick, PastelBlue2, fill=PastelBlue2, fill opacity=0.25] (1.4,0) rectangle (3.6,2.0);
            \node[PastelBlue2, font=\scriptsize] at (3.4,-0.3) {$B$};
            \draw[ultra thick, dashed, black] (1.4,0.6) rectangle (2.6,2.0);
        \end{tikzpicture}
    \end{column}
\end{columns}
\end{frame}

\begin{frame}{IoU vs Dice: Exemplo Numérico}
\begin{itemize}
    \item Suponha duas caixas com:
    \begin{itemize}
        \item $\mathrm{Area}(A) = 10$.
        \item $\mathrm{Area}(B) = 10$.
        \item $A \cap B = 5$.
        \item $A \cup B = \mathrm{Area}(A) + \mathrm{Area}(B) - A \cap B = 15$.
    \end{itemize}
\end{itemize}
\vspace{0.2em}
\begin{columns}
    \begin{column}{.5\linewidth}
        \begin{equation*}
            \mathrm{IoU} = \frac{5}{15} \approx 0.333
        \end{equation*}
    \end{column}
    \begin{column}{.5\linewidth}
        \begin{equation*}
            \mathrm{Dice} = \frac{2 \cdot 5}{10 + 10} = 0.5
        \end{equation*}
    \end{column}
\end{columns}
\vspace{0.6em}
\begin{itemize}
    \item Ambas as métricas medem a sobreposição entre caixas.
    \item Para a mesma sobreposição, o IoU produz valores menores que o \textit{Dice}.
    \begin{itemize}
        \item Pode ser interpretado como uma penalização ``mais quadrática'' do erro.
    \end{itemize}
\end{itemize}
\end{frame}

\begin{frame}{Localização: Resumo}
\begin{itemize}
    \item A \textit{bounding box} é modelada como quatro saídas de regressão multivariada.
    \item Função de custo: MSE.
    \item Ativação na saída: sigmoide (4 unidades).
    \item Podemos:
    \begin{itemize}
        \item Reaproveitar pesos de convolução já treinados.
        \item Treinar apenas a nova cabeça densa.
        \item Ou treinar a rede inteira (\textit{fine-tuning}).
    \end{itemize}
    \item Avaliamos com IoU ou \textit{Dice Score}, comparando com o \textit{ground truth}.
\end{itemize}
\end{frame}

%========================================================================
\section{Classificação e Localização}
%========================================================================

\begin{frame}{Múltiplas Tarefas Simultâneas}
\begin{itemize}
    \item Podemos resolver classificação e localização \textbf{ao mesmo tempo}.
    \item A parte convolucional serve como \textbf{extrator de \textit{features}} compartilhado.
    \item Adicionamos duas cabeças densas independentes:
    \begin{itemize}
        \item Cabeça de classificação: \textit{softmax} sobre as 1000 classes.
        \item Cabeça de localização: 4 unidades sigmoide para a \textit{bounding box}.
    \end{itemize}
    \item Cada cabeça contribui com seu próprio termo na função de custo:
\end{itemize}
\begin{equation*}
\mathcal{L} = \mathcal{L}_{\text{classif.}} + \lambda\,\mathcal{L}_{\text{loc.}}
\end{equation*}
\end{frame}

\begin{frame}{Backbone como Extrator de Features}
\begin{columns}[c]
    \begin{column}{.42\linewidth}
        \begin{itemize}
            \item O \textit{backbone} convolucional é \textbf{compartilhado} pelas duas cabeças.
            \item Cada cabeça densa recebe as mesmas \textit{features} e produz sua própria saída.
            \item Trocando a cabeça e mantendo o \textit{backbone}, podemos resolver muitas tarefas distintas a partir das mesmas \textit{features}.
        \end{itemize}
    \end{column}
    \begin{column}{.58\linewidth}
        \centering
        \renewcommand{\vgpitch}{0.235}%
        \begin{tikzpicture}[every node/.style={font=\tiny}]
            % backbone compartilhado (entrada + convoluções + pooling)
            \foreach \i/\txt/\col/\tcol in {
                0/entrada/black!55/white,
                1/\chipconv{64}/BrickRed!75/white,
                2/\chipconv{64}/BrickRed!75/white,
                3/pooling/PastelGreen/black,
                4/\chipconv{128}/BrickRed!75/white,
                5/\chipconv{128}/BrickRed!75/white,
                6/pooling/PastelGreen/black,
                7/\chipconv{256}/BrickRed!75/white,
                8/\chipconv{256}/BrickRed!75/white,
                9/\chipconv{256}/BrickRed!75/white,
                10/pooling/PastelGreen/black,
                11/\chipconv{512}/BrickRed!75/white,
                12/\chipconv{512}/BrickRed!75/white,
                13/\chipconv{512}/BrickRed!75/white,
                14/pooling/PastelGreen/black,
                15/\chipconv{512}/BrickRed!75/white,
                16/\chipconv{512}/BrickRed!75/white,
                17/\chipconv{512}/BrickRed!75/white,
                18/pooling/PastelGreen/black}{
                \node[vgchipbase, fill=\col, text=\tcol, draw=black!50, thin] at (0, \i*\vgpitch) {\txt};
            }
            \draw[decorate, decoration={brace, amplitude=6pt, mirror}, BrickRed, thick]
                (1.05, 1*\vgpitch - 0.1) -- (1.05, 18*\vgpitch + 0.1)
                node[midway, right=8pt, BrickRed, font=\scriptsize\bfseries, align=center] {\textit{backbone}\\VGG\\(compartilhado)};

            % imagem de entrada
            \node[inner sep=0, anchor=east] (img) at (-1.3, 0)
              {\includegraphics[width=1.5cm]{slides/img/02-cnn/loc_cat_snow.jpg}};
            \node[anchor=north] at (img.south) {Imagem};
            \draw[->, thick] (img.east) -- (-0.98, 0);

            % duas cabeças densas
            \pgfmathsetmacro{\hy}{18*\vgpitch + 0.72}
            \foreach \hx/\last/\act/\actcol/\acttext in {-1.05/\chipfc{1000}/softmax/Goldenrod!70/black, 1.05/\chipfc{4}/sigmoide/Goldenrod!70/black}{
                \node[vgchipbase, fill=PastelBlue2!80, text=white, draw=black!50, thin] at (\hx, \hy) {\chipfc{4096}};
                \node[vgchipbase, fill=PastelBlue2!80, text=white, draw=black!50, thin] at (\hx, \hy + \vgpitch) {\chipfc{4096}};
                \node[vgchipbase, fill=PastelBlue2!80, text=white, draw=black!50, thin] at (\hx, \hy + 2*\vgpitch) {\last};
                \node[vgchipbase, fill=\actcol, text=\acttext, draw=black!50, thin] at (\hx, \hy + 3*\vgpitch) {\act};
            }
            % bifurcação
            \draw[thick] (0, 18*\vgpitch + 0.11) -- (0, \hy - 0.38);
            \draw[->, thick] (0, \hy - 0.38) -| (-1.05, \hy - 0.12);
            \draw[->, thick] (0, \hy - 0.38) -| (1.05, \hy - 0.12);
            % saídas
            \node[anchor=south, align=center, font=\scriptsize] at (-1.05, \hy + 3*\vgpitch + 0.12) {\textbf{classe}};
            \node[anchor=south, align=center, font=\scriptsize] at (1.05, \hy + 3*\vgpitch + 0.12) {\textbf{caixa} $(x,y,h,w)$};
        \end{tikzpicture}
    \end{column}
\end{columns}
\end{frame}

%========================================================================
\section{Segmentação Semântica}
%========================================================================

\begin{frame}{Segmentação Semântica}
\begin{itemize}
    \item É uma classificação em \textbf{nível de pixel}.
    \item Cada pixel da imagem recebe uma probabilidade de pertencer a cada classe.
    \item Aplicações típicas: direção autônoma, imagens médicas, sensoriamento remoto.
\end{itemize}
\vspace{0.5em}
\centering
\begin{tikzpicture}[every node/.style={font=\scriptsize}]
    \node[inner sep=0] (in)
      {\includegraphics[width=4.7cm]{slides/img/02-cnn/segmentation_example.jpg}};
    \node[above=2pt of in] {Imagem de entrada};
    \node[inner sep=0, right=1.3cm of in] (out)
      {\includegraphics[width=4.7cm]{slides/img/02-cnn/segmentation_example_mask.png}};
    \node[above=2pt of out] {Máscara por classe};
    \draw[->, very thick] (in.east) -- (out.west);
\end{tikzpicture}\\[4pt]
{\tiny Cada pixel recebe uma classe (\textit{notebook}, copo, mesa, fundo). Exemplo CC0 de Martin Thoma, \textit{Wikimedia Commons}.}
\end{frame}

\begin{frame}{Primeira Ideia (Ingênua)}
\begin{itemize}
    \item Separar a imagem em \textit{patches}.
    \item Para cada \textit{patch}, classificar o pixel central com uma CNN.
    \item Repetir para todos os pixels da imagem.
    \item Problemas:
    \begin{itemize}
        \item Extremamente \textbf{ineficiente}: requer milhares de \textit{forwards} por imagem.
        \item Muito redundante, pois \textit{patches} vizinhos compartilham a maior parte do conteúdo.
    \end{itemize}
    \item Precisamos de uma forma de classificar todos os pixels em um único \textit{forward}.
\end{itemize}
\end{frame}

\begin{frame}{Arquitetura Encoder-Decoder}
\begin{itemize}
    \item A CNN deve classificar todos os pixels em um único \textit{forward}\footfullcite{noh2015deconv}.
    \item \textbf{Encoder}: reduz a dimensão espacial via convoluções e \textit{pooling}.
    \item \textbf{Decoder}: aumenta a dimensão espacial via \textit{unpooling} e convoluções transpostas.
\end{itemize}
\vspace{0.2em}
\centering
\begin{tikzpicture}[every node/.style={font=\tiny}]
    \def\bw{0.35}

    % Encoder pyramid (decreasing)
    \draw[fill=BrickRed!30] (0,0.0) rectangle ++(\bw,1.8);
    \draw[fill=BrickRed!40] (0.6,0.2) rectangle ++(\bw,1.4);
    \draw[fill=BrickRed!55] (1.2,0.4) rectangle ++(\bw,1.0);
    \draw[fill=BrickRed!70] (1.8,0.55) rectangle ++(\bw,0.7);
    \draw[fill=BrickRed!85] (2.4,0.7) rectangle ++(\bw,0.4);

    % Bottleneck
    \draw[fill=black!50] (3.1,0.8) rectangle ++(0.25,0.2);
    \node at (3.22,1.25) {bottleneck};

    % Decoder (increasing)
    \draw[fill=PastelBlue2, opacity=0.35] (3.6,0.7) rectangle ++(\bw,0.4);
    \draw[fill=PastelBlue2, opacity=0.5] (4.2,0.55) rectangle ++(\bw,0.7);
    \draw[fill=PastelBlue2, opacity=0.6] (4.8,0.4) rectangle ++(\bw,1.0);
    \draw[fill=PastelBlue2, opacity=0.75] (5.4,0.2) rectangle ++(\bw,1.4);
    \draw[fill=PastelBlue2, opacity=0.9] (6.0,0.0) rectangle ++(\bw,1.8);

    % Labels
    \node[align=center] at (1.2,-0.3) {\scriptsize Encoder (VGG16)};
    \node[align=center] at (4.8,-0.3) {\scriptsize Decoder (conv transposta)};

    % Arrows
    \draw[->, thick, dashed] (-0.2,2.0) -- (3.0,2.0) node[midway, above] {redução espacial};
    \draw[->, thick, dashed] (3.4,2.0) -- (6.6,2.0) node[midway, above] {expansão espacial};
\end{tikzpicture}

\vspace{0.2em}
\begin{itemize}
    \item Saída: máscara com a mesma resolução da entrada, com uma classe por pixel.
\end{itemize}
\end{frame}

%========================================================================
\section{Detecção de Objetos}
%========================================================================

\begin{frame}{Detecção de Objetos}
\begin{itemize}
    \item Tarefa muito parecida com classificação $+$ localização.
    \item Diferença: a imagem pode conter \textbf{vários objetos} simultaneamente.
    \begin{itemize}
        \item Cada objeto detectado tem sua própria classe e \textit{bounding box}.
    \end{itemize}
    \item Desafio principal: \textbf{não sabemos de antemão quantos objetos} aparecem em cada imagem.
\end{itemize}
\vspace{0.1em}
\centering
\resizebox{5.1cm}{!}{%
\begin{tikzpicture}[every node/.style={font=\scriptsize\bfseries}]
    \node[inner sep=0, anchor=south west] (img) at (0,0)
      {\includegraphics[width=6.5cm]{slides/img/02-cnn/det_cat_dog.jpg}};
    % cachorro (esquerda)
    \draw[BrickRed, very thick] (0.45,2.07) rectangle (5.11,4.86);
    \node[BrickRed, anchor=south west, fill=white, inner sep=1pt] at (0.45,4.86) {cachorro};
    % gato (direita)
    \draw[PastelBlue2, very thick] (2.80,0.54) rectangle (6.08,2.72);
    \node[PastelBlue2, anchor=south east, fill=white, inner sep=1pt] at (6.08,2.72) {gato};
\end{tikzpicture}}\\[3pt]
{\tiny Cada objeto recebe classe $+$ \textit{bounding box}. Foto: \textit{Socks e Buddy} (Casa Branca), domínio público.}
\end{frame}

\begin{frame}{Famílias de Detectores: Two-Stage}
\begin{itemize}
    \item \textbf{R-CNN}\footfullcite{girshick2014rcnn}: gera $\sim 2.000$ \textit{region proposals} (via \textit{selective search}), recorta cada uma e roda a CNN separadamente. Lento, mas estabeleceu a base do paradigma.
    \item \textbf{Fast R-CNN}: a CNN processa a imagem inteira uma só vez; \textit{ROI pooling} extrai \textit{features} de cada região a partir do mesmo \textit{feature map}.
    \item \textbf{Faster R-CNN}\footfullcite{ren2015fasterrcnn}: substitui \textit{selective search} por uma \textbf{Region Proposal Network} (RPN), também convolucional. Treino \textit{end-to-end}.
    \item Característica comum: primeiro propõem regiões, depois classificam cada uma. Tipicamente mais precisos, porém mais lentos.
\end{itemize}
\end{frame}

\begin{frame}{Famílias de Detectores: One-Stage e Transformers}
\begin{itemize}
    \item \textbf{YOLO}\footfullcite{redmon2016yolo} (\textit{You Only Look Once}): divide a imagem em uma grade $S\times S$ e prediz, em \textbf{uma única passada}, a classe e a \textit{bounding box} para cada célula.
    \begin{itemize}
        \item Muito mais rápido (opera em tempo real), ao custo de menor precisão em objetos pequenos.
        \item Variantes posteriores (YOLOv3-v8) recuperam grande parte dessa diferença de precisão.
    \end{itemize}
    \item \textbf{SSD} (\textit{Single Shot MultiBox Detector}): também \textit{one-stage}, prediz caixas em vários níveis de resolução simultaneamente.
    \item \textbf{DETR}\footfullcite{carion2020detr}: formula detecção como predição direta de um conjunto de caixas usando um \textbf{Transformer}.
    \begin{itemize}
        \item Elimina \textit{anchors} e \textit{non-maximum suppression}.
        \item Treina com \textit{Hungarian matching} entre predições e \textit{ground truth}.
    \end{itemize}
\end{itemize}
\end{frame}

%========================================================================
\section{Segmentação de Instâncias}
%========================================================================

\begin{frame}{Segmentação de Instâncias}
\begin{itemize}
    \item Muito parecida com detecção de objetos:
    \begin{itemize}
        \item Cada objeto recebe uma máscara em nível de pixel, além da \textit{bounding box}.
        \item Instâncias diferentes da mesma classe são \textbf{diferenciadas} entre si.
    \end{itemize}
    \item Diferença em relação à segmentação semântica: dois gatos lado a lado aparecem como duas máscaras distintas, não como uma só.
\end{itemize}
\end{frame}

\begin{frame}{Segmentação Semântica vs Instâncias}
\centering
\begin{tikzpicture}[every node/.style={font=\scriptsize}]
    \node[inner sep=0] (sem)
      {\includegraphics[width=5.5cm]{slides/img/02-cnn/seg_semantic.png}};
    \node[above=2pt of sem] {\textbf{Segmentação Semântica}};
    \node[inner sep=0, right=0.8cm of sem] (inst)
      {\includegraphics[width=5.5cm]{slides/img/02-cnn/seg_instance.png}};
    \node[above=2pt of inst] {\textbf{Segmentação de Instâncias}};
\end{tikzpicture}\\[4pt]
{\small Na \textbf{semântica}, todos os objetos da classe ``balão'' recebem a \textbf{mesma} cor;
na de \textbf{instâncias}, cada balão recebe uma cor \textbf{distinta}.}\\[2pt]
{\tiny Figura: \textit{Mask R-CNN}, \textit{Wikimedia Commons} (CC BY-SA 4.0).}
\end{frame}

%========================================================================
\section{Referências}
%========================================================================

\begin{frame}{Leituras Recomendadas}
    \leiturasize
\begin{itemize}
    \item Sugere-se \textbf{fortemente} a leitura do Capítulo 10 de \textit{Understanding Deep Learning}~\cite{prince2023understanding}.
    \begin{itemize}
        \item Disponível em \url{https://udlbook.github.io/udlbook/}.
    \end{itemize}
    \item Os artigos abaixo são acessíveis e fornecem a visão original dos autores sobre as ideias apresentadas nesta aula:
    \begin{itemize}
        \item \fullcite{simonyan2014vgg}
        \item \fullcite{sermanet2013overfeat}
        \item \fullcite{noh2015deconv}
    \end{itemize}
\end{itemize}
\end{frame}

\begin{frame}{Referências}
\printbibliography
\end{frame}

\end{document}
